\section{Method}
\label{sec:method}

\subsection{Trajectory Preprocessing and Window Construction}
% Freeze coordinate origin and units, resampling, scaling, missing-data
% policy, minimum usable fragment criteria, and duplicate-flight handling.
% Record one aircraft-window per output embedding.
\draftnote{Describe the versioned transformation from observations to nine-dimensional states.}

\subsection{Self-Supervised Trajectory Encoder}
% Describe RDP segment identities, positive/negative pairs, causal
% Transformer architecture, masking, and the exact contrastive loss.
% Record the difference from ATSCC/ITSC configurations.
% Freeze the selected encoder before comparing scenes.
\draftnote{Specify the training objective, architecture, and checkpoint-selection protocol.}

\subsection{Temporal Pooling and Distance Geometry}
% Define sum, mean, and signed GeM pooling, including padding masks.
% Record whether pooled embeddings are L2-normalized.
% Under final per-vector L2 normalization, sum and mean are identical
% for the same nonzero hidden-state sum.
\draftnote{Define the pooling and distance choices used by each experiment.}

\subsection{Empirical Correlation Dimension}
For a fixed distance function $d(\cdot,\cdot)$ and $N_s\geq2$, the
empirical correlation sum is
\begin{equation}
\Corr_s(r)=
\frac{2}{N_s(N_s-1)}
\sum_{i<j}\indic{d(z_{i,s},z_{j,s})<r}.
\label{eq:correlation-sum}
\end{equation}
Following the correlation-dimension construction
\cite{grassberger1983measuring}, an empirical slope can be estimated over
a prespecified admissible scale interval:
\begin{equation}
\log\Corr_s(r)\approx \Dscene(s)\log r+b_s.
\label{eq:empirical-slope}
\end{equation}
This finite-sample slope is a candidate geometric descriptor. Its
association with operational complexity remains an empirical question.
\draftnote{Specify scale selection, sample-size checks, slope stability,
normalization, and conditions under which a scene is unestimable.}
